% !PRINTED 756
\weberpage{756}
%<<CONT
dass $C_0 F(0) + C_1 F(1) + \cdots + C_m F(m)$ eine nicht durch $p$ theilbare, also auch nicht verschwindende ganze Zahl ist, und damit ist 1) bewiesen.

Wenn wir nun noch nachweisen können, dass $Q(x)$ für jedes positive $x$ durch Vergrösserung von $p$ beliebig klein gemacht werden kann, so folgt, dass bei hinlänglich grossem $p$ der Ausdruck 2) gewiss kleiner als 1 ist, und unser Beweis ist vollendet.

Gehen wir aber zu dem Ausdrucke (10) für $Q(x)$ zurück und bezeichnen mit $\gamma_n, \gamma_{n-1}, \ldots, \gamma_0$ die absoluten Werthe der Coëfficienten $c_n, c_{n-1}, \ldots, c_0$ von $f(x)$, so sehen wir, da die $q_n, q_{n-1}, \ldots$ dem absoluten Werthe nach kleiner als 1 sind, dass für jedes positive $x$ der absolute Werth von $Q(x)$ kleiner ist als

\weq{17}{\psi(x) = \gamma_n x^n + \gamma_{n-1}x^{n-1}+\cdots+\gamma_0.}

Die Coëfficienten $c_n, c_{n-1}, \ldots, c_0$ der Function $f(x)$ in (15) unterscheiden sich aber von den $\gamma_n, \gamma_{n-1}, \ldots, \gamma_0$ nur durch das Vorzeichen. Ersetzen wir daher $x$ durch $-x$, bilden also die Function
\[
f(-x) = \frac{x^{p-1}(1+x)^p(2+x)^p\ldots(m+x)^p}{\Pi(p-1)},
\]
so hat diese Function dieselben Coëfficienten, wie $f(x)$, aber alle mit positiven Vorzeichen. Sie ist also keine andere, als die Function $\psi(x)$.

Setzen wir also noch
\[
X = x(1+x)(2+x)\ldots(m+x),
\]
so wird

\weq{18}{\psi(x) = \frac{X}{x}\cdot\frac{X^{p-1}}{\Pi(p-1)},}

und dies nähert sich, wie am Anfange dieses Paragraphen nachgewiesen ist, mit unendlich wachsendem $p$ der Grenze \gesperrt{Null}.

\gesperrt{Es ist also $e$ eine transcendente Zahl.}

\begin{center}
§. 206.

\gesperrt{Transcendenz der Zahl $\pi$.}
\end{center}

Mit denselben Hülfsmitteln lässt sich nun auch die Transcendenz der Zahl $\pi$ beweisen. Als Definition dieser Zahl dient uns dabei, dass es die kleinste positive Zahl ist, die der Gleichung

\weq{1}{e^{i\pi} = -1}

genügt, wenn $e^x$ durch die Formel \S.\,205, (2) definirt wird.
